\documentclass[14pt,aspectratio=1610]{beamer}
% \documentclass[12pt,aspectratio=1610,handout]{beamer}
\usepackage{csu-theme}
\usepackage{arrays}
\usepackage{linked-lists}
\usepackage{booktabs}
\usetikzlibrary{backgrounds}

\title{\Huge{} Standard Template Library\\Algorithms}
\subtitle{CSCI 315: \textit{Data Structure Analysis}\\Chapter 21}
%\author{Sean T. Hayes, Ph.D.}
\institute{Charleston Southern University}

%% Comment out date to use default (today)
% \date{March 20, 2025}
\subject{Software Programming}
%\keywords{}
\graphicspath{{../imgs/stl}}

\setlength{\parskip}{0.5em}

\begin{document}

\begin{frame}
  \titlepage{}
\end{frame}

\section*{Algorithms}

\begin{frame}{Algorithms in the STL}
Many algorithms that work on any STL container are \\
found in the \lstinline{<algorithm>} library.
\begin{itemize}
\item
  Copy
\item
  Partition
\item
  Sort
\item
  Search
\item
  \ldots{}
\end{itemize}

These algorithms are provided generically to support many container types.
\end{frame}


\begin{frame}{STL Algorithm Classifications}

\begin{itemize}
\item
  \emph{Non-modifying algorithms} do not modify the elements of the container
\item
  \emph{Modifying algorithms} modify the elements of the container
\item
  \emph{Numeric algorithms} perform numeric calculations on elements of a container.
\item
  \emph{Heap algorithms} are Heapsort algorithm that sorts contiguous data.
\end{itemize}
\end{frame}

\begin{frame}{Non-Modifying Sequence Algorithms}
\footnotesize\renewcommand{\ULdepth}{3pt}\renewcommand{\arraystretch}{1.2}
\begin{tabularx}{\textwidth}{r X}

\href{https://en.cppreference.com/w/cpp/algorithm/all_any_none_of}{\lstinline{all\_of, any\_of, none\_of}}
 &    checks if a predicate is true for all, any or none of the elements in a range \\
\href{https://en.cppreference.com/w/cpp/algorithm/for_each}{\lstinline{for\_each}}
 &    applies a function to a range of elements \\
\href{https://en.cppreference.com/w/cpp/algorithm/for_each_n}{\lstinline{for\_each\_n}}
 &    applies a function object to the first n elements of a sequence  \\
\href{https://en.cppreference.com/w/cpp/algorithm/count}{\lstinline{count, count\_if}}
 &    returns the number of elements satisfying specific criteria  \\
\href{https://en.cppreference.com/w/cpp/algorithm/mismatch}{\lstinline{mismatch}}
 &    finds the first position where two ranges differ \\
\href{https://en.cppreference.com/w/cpp/algorithm/find}{\lstinline{find, find\_if, find\_if\_not}}
 &    finds the first element satisfying specific criteria  \\
\href{https://en.cppreference.com/w/cpp/algorithm/find_end}{\lstinline{find\_end}}
 &    finds the last sequence of elements in a certain range \\
\href{https://en.cppreference.com/w/cpp/algorithm/find_first_of}{\lstinline{find\_first\_of}}
 &    searches for any one of a set of elements  \\
\href{https://en.cppreference.com/w/cpp/algorithm/adjacent_find}{\lstinline{adjacent\_find}}
 &    finds the first two adjacent items that are equal (or satisfy a given predicate)  \\
\href{https://en.cppreference.com/w/cpp/algorithm/search}{\lstinline{search}}
 &    searches for a range of elements \\
\href{https://en.cppreference.com/w/cpp/algorithm/search_n}{\lstinline{search\_n}}
 &    searches a range for a number of consecutive\\[-0.2em]
 &    copies of an element \\
\end{tabularx}
\end{frame}


\begin{frame}[fragile]{Example: The \texttt{for\_each()} Algorithm}
  \small
\begin{lstlisting}
#include <vector>
#include <algorithm>
#include <iostream>
void show(int n) {
  std::cout << n << ", ";
}

int main() {
  // initialize vector with a list of values
  std::vector<int> v{ 12, 3, 17, 8 };

  // apply function show to each element of vector v
  std::for_each (v.begin(), v.end(), show);
}
\end{lstlisting}
\end{frame}

\begin{frame}[fragile]{Example: The \texttt{find()} Algorithm}
  
\begin{lstlisting}[basicstyle=\small\ttfamily]
#include <vector>
#include <algorithm>
#include <iostream>
int main() {
  std::vector<int> v{ 12, 3, 17, 8, 34, 56, 9 };
  std::vector<int>::iterator iter;
  int key;

  std::cout << "Enter value: ";
  std::cin >> key;
  iter = std::find(v.begin(), v.end(), key);
  std::cout << "Element " << key << " is ";
  if (iter == v.end()) // value was NOT found
     std::cout << "not";
  std::cout << " in vector v.\n";
}
\end{lstlisting}
\end{frame}

\begin{frame}[fragile]{Example: The \texttt{find\_if()} Algorithm}
\begin{lstlisting}[basicstyle=\small\ttfamily]
#include <vector>
#include <algorithm>
#include <iostream>
bool inRange(int n) { return (n > 21) && (n < 36); };

int main() {
  std::vector<int> v{ 12, 3, 17, 8, 34, 56, 9 };
  std::vector<int>::iterator iter;
  iter = std::find_if(v.begin(), v.end(), inRange);

  // finds an element in v for which inRange is true
  if (iter != v.end()) // found the element
    std::cout << "found " << *iter << '\n';
  else
    std::cout << "not found\n";
}
\end{lstlisting}
\end{frame}

\begin{frame}[fragile]{Example: The \texttt{count\_if()} Algorithm}
\begin{lstlisting}[basicstyle=\small\ttfamily]
#include <vector>
#include <algorithm>
#include <iostream>
bool inRange(int n) { return (n > 14) && (n < 36); };

int main() {
  std::vector<int> v{ 12, 3, 17, 8, 34, 56, 9 };
  std::vector<int>::iterator iter;

  // counts element in v for which inRange is true
  const int N = std::count_if(v.begin(),v.end(), inRange);
  std::cout << "found " << N << " values between 14 and 36.\n";
}
\end{lstlisting}
\end{frame}

\begin{frame}{Minimum/Maximum Algorithms}
\small\renewcommand{\ULdepth}{3pt}\renewcommand{\arraystretch}{1.3}
\begin{tabularx}{\textwidth}{r X}
\href{https://en.cppreference.com/w/cpp/algorithm/max}{\lstinline{max}} &
  returns the greater of the given values\\
\href{https://en.cppreference.com/w/cpp/algorithm/max_element}{\lstinline{max_element}} &
  returns the largest element in a range\\
\href{https://en.cppreference.com/w/cpp/algorithm/min}{\lstinline{min}} &
  returns the smaller of the given values\\
\href{https://en.cppreference.com/w/cpp/algorithm/min_element}{\lstinline{min_element}} &
  returns the smallest element in a range\\
\href{https://en.cppreference.com/w/cpp/algorithm/minmax}{\lstinline{minmax}} &
  returns the smaller and larger of two elements\\
\href{https://en.cppreference.com/w/cpp/algorithm/minmax_element}{\lstinline{minmax_element}} &
  returns the smallest and the largest elements in a range\\
\href{https://en.cppreference.com/w/cpp/algorithm/clamp}{\lstinline{clamp}} &
  clamps a value between a pair of boundary values\\
\end{tabularx}
\end{frame}

\begin{frame}{Set Algorithms (on Sorted Ranges)}
\small\renewcommand{\ULdepth}{3pt}\renewcommand{\arraystretch}{1.3}
\begin{tabularx}{\textwidth}{r X}
\href{https://en.cppreference.com/w/cpp/algorithm/includes}{\lstinline{includes}} &
  returns true if one sequence is a subsequence of another\\
\href{https://en.cppreference.com/w/cpp/algorithm/set_difference}{\lstinline{set_difference}} &
  computes the difference between two sets\\
\href{https://en.cppreference.com/w/cpp/algorithm/set_intersection}{\lstinline{set_intersection}} &
  computes the intersection of two sets\\
\href{https://en.cppreference.com/w/cpp/algorithm/set_symmetric_difference}{\lstinline{set_symmetric_difference}} &
  computes the symmetric difference between two sets\\
\href{https://en.cppreference.com/w/cpp/algorithm/set_union}{\lstinline{set_union}} &
  computes the union of two sets\\
\end{tabularx}
\end{frame}

\begin{frame}{Modifying Sequence Algorithms}
\footnotesize\renewcommand{\ULdepth}{3pt}\renewcommand{\arraystretch}{1.3}
\begin{tabularx}{\textwidth}{r X}
\href{https://en.cppreference.com/w/cpp/algorithm/copy}{\lstinline{copy, copy_if}} &
  copies a range of elements to a new location\\
\href{https://en.cppreference.com/w/cpp/algorithm/copy_n}{\lstinline{copy_n}} &
  copies a number of elements to a new location\\
\href{https://en.cppreference.com/w/cpp/algorithm/copy_backward}{\lstinline{copy_backward}} &
  copies a range of elements in backward order\\
\href{https://en.cppreference.com/w/cpp/algorithm/move}{\lstinline{move}} &
  moves a range of elements to a new location\\
\href{https://en.cppreference.com/w/cpp/algorithm/move_backward}{\lstinline{move_backward}} &
  moves a range of elements to a new location in backward order\\
\href{https://en.cppreference.com/w/cpp/algorithm/fill}{\lstinline{fill}} &
  copy-assigns the given value to every element in a range\\
\href{https://en.cppreference.com/w/cpp/algorithm/fill_n}{\lstinline{fill_n}} &
  copy-assigns the given value to N elements in a range\\
\href{https://en.cppreference.com/w/cpp/algorithm/transform}{\lstinline{transform}} &
  applies a function to a range of elements, storing results in a destination range\\
\href{https://en.cppreference.com/w/cpp/algorithm/generate}{\lstinline{generate}} &
  assigns the results of successive function calls to every element in a range\\
\href{https://en.cppreference.com/w/cpp/algorithm/generate_n}{\lstinline{generate_n}} &
  assigns the results of successive function calls to N elements in a range\\
\href{https://en.cppreference.com/w/cpp/algorithm/remove}{\lstinline{remove, remove_if}} &
  removes elements satisfying specific criteria\\

\end{tabularx}
\end{frame}

\begin{frame}[fragile]{Example: The \texttt{copy()} Algorithm}

The copy algorithm overwrites the existing contents when using regular iterators.

\begin{lstlisting}[basicstyle=\small\ttfamily]
#include <list>
#include <algorithm>

int main() {
  std::list<int> odds{ 1, 3, 5, 7, 9 };
  std::list<int> result{ 2, 4, 6, 8, 10 };

  // copy contents of odds to result overwriting the elements in result
  std::copy(odds.begin(), odds.end(), result.begin());
  // result = { 1, 3, 5, 7, 9 }
}
\end{lstlisting}
\end{frame}

\begin{frame}[fragile]{Example: The \texttt{copy()} Algorithm}
\small
With insert operators, you can modify the behavior of the copy algorithm.\\[0.5em]

\begin{tabular}{l l}
\href{https://en.cppreference.com/w/cpp/iterator/back_inserter}{\lstinline{back_inserter}} & inserts new elements at the end\\
\href{https://en.cppreference.com/w/cpp/iterator/front_inserter}{\lstinline{front_inserter}} & inserts new elements at the beginning\\
\href{https://en.cppreference.com/w/cpp/iterator/inserter}{\lstinline{inserter}} & inserts new elements at a specified location\\
\end{tabular}\\[1em]

\begin{lstlisting}[basicstyle=\small\ttfamily]
std::list<int> l1{ 1, 3, 5, 7, 9 };
std::list<int> l2{ 2, 4, 6, 8, 10 };
\end{lstlisting}

\begin{lstlisting}[basicstyle=\footnotesize\ttfamily]
std::copy(l1.begin(), l1.end(), std::back_inserter(l2));
// l2 = { 2, 4, 6, 8, 10, 1, 3, 5, 7, 9  }
\end{lstlisting}

\begin{lstlisting}[basicstyle=\footnotesize\ttfamily]
std::copy(l1.begin(), l1.end(), std::front_inserter(l2));
// l2 = { 9, 7, 5, 3, 1, 2, 4, 6, 8, 10 }
\end{lstlisting}

\begin{lstlisting}[basicstyle=\footnotesize\ttfamily]
std::copy(l1.begin(), l1.end(), std::inserter(l2, l2.begin()));
// l2 = { 1, 3, 5, 7, 9, 2, 4, 6, 8, 10 }
\end{lstlisting}
\end{frame}

\begin{frame}[fragile]{Example: \texttt{copy()} with Output Iterator}

\begin{lstlisting}[basicstyle=\small\ttfamily]
#include <iostream>
#include <iterator>
#include <vector>
#include <algorithm>

int main() {
  std::vector<int> vector {10, 20, 30, 40, 50, 60};
  std::ostream_iterator<int> outIter{std::cout,
    ", "};

  std::copy(vector.begin(), vector.end(), outIter);
}
\end{lstlisting}
\end{frame}

\begin{frame}{Modifying Sequence Algorithms}
\footnotesize\renewcommand{\ULdepth}{3pt}\renewcommand{\arraystretch}{1.3}
\begin{tabularx}{\textwidth}{r X}
\href{https://en.cppreference.com/w/cpp/algorithm/remove_copy}{\lstinline{remove_copy, remove_copy_if}} &
  copies a range of elements omitting those that satisfy specific criteria\\
\href{https://en.cppreference.com/w/cpp/algorithm/replace}{\lstinline{replace, replace_if}} &
  replaces all values satisfying specific criteria with another value\\
\href{https://en.cppreference.com/w/cpp/algorithm/replace_copy}{\lstinline{replace_copy, replace_copy_if}} &
  copies a range, replacing elements satisfying specific criteria with another value\\
\href{https://en.cppreference.com/w/cpp/algorithm/swap}{\lstinline{swap}} &
  swaps the values of two objects\\
\href{https://en.cppreference.com/w/cpp/algorithm/swap_ranges}{\lstinline{swap_ranges}} &
  swaps two ranges of elements\\
\href{https://en.cppreference.com/w/cpp/algorithm/iter_swap}{\lstinline{iter_swap}} &
  swaps the elements pointed to by two iterators\\
\href{https://en.cppreference.com/w/cpp/algorithm/reverse}{\lstinline{reverse}} &
  reverses the order of elements in a range\\
\href{https://en.cppreference.com/w/cpp/algorithm/reverse_copy}{\lstinline{reverse_copy}} &
  creates a copy of a range that is reversed\\
\href{https://en.cppreference.com/w/cpp/algorithm/rotate}{\lstinline{rotate}} &
  rotates the order of elements in a range\\
\href{https://en.cppreference.com/w/cpp/algorithm/rotate_copy}{\lstinline{rotate_copy}} &
  copies and rotate a range of elements\\
\end{tabularx}
\end{frame}

\begin{frame}[fragile]{Example: The \texttt{sort()} and \texttt{merge()} Algorithms}

\begin{lstlisting}[basicstyle=\small\ttfamily]
#include <list>
#include <algorithm>
int main() {
  std::list<int>  l1 { 6, 4, 9, 1, 7 }
  std::list<int>  l2 { 4, 2, 1, 3, 8 };
  l1.sort(); // l1 = {1, 4, 6, 7, 9}
  l2.sort(); // l2 = {1, 2, 3, 4, 8 }
  l1.merge(l2);  // merges l2 into l1
  // l1 = { 1, 1, 2, 3, 4, 4, 6, 7, 8, 9},  l2 = {}
}
\end{lstlisting}
\end{frame}

\section{Numeric Algorithms}

\begin{frame}{Numeric Algorithms in \href{https://en.cppreference.com/w/cpp/header/numeric}{\texttt{<numeric>}}}
\footnotesize\renewcommand{\ULdepth}{3pt}\renewcommand{\arraystretch}{1.3}
\begin{tabularx}{\textwidth}{r X}
\href{https://en.cppreference.com/w/cpp/algorithm/iota}{\lstinline{iota}} &
  fills a range with successive increments of the starting value\\
\href{https://en.cppreference.com/w/cpp/algorithm/accumulate}{\lstinline{accumulate}} &
  sums up a range of elements\\
\href{https://en.cppreference.com/w/cpp/algorithm/inner_product}{\lstinline{inner_product}} &
  computes the inner product of two ranges of elements\\
\href{https://en.cppreference.com/w/cpp/algorithm/adjacent_difference}{\lstinline{adjacent_difference}} &
  computes the differences between adjacent elements in a range\\
\href{https://en.cppreference.com/w/cpp/algorithm/partial_sum}{\lstinline{partial_sum}} &
  computes the partial sum of a range of elements\\
\href{https://en.cppreference.com/w/cpp/algorithm/reduce}{\lstinline{reduce}} &
  similar to \texttt{std::accumulate}, except out of order\\
\href{https://en.cppreference.com/w/cpp/algorithm/exclusive_scan}{\lstinline{exclusive_scan}} &
  excludes the ith input element from the ith sum\\
\href{https://en.cppreference.com/w/cpp/algorithm/inclusive_scan}{\lstinline{inclusive_scan}} &
  includes the ith input element in the ith sum\\
\href{https://en.cppreference.com/w/cpp/algorithm/transform_reduce}{\lstinline{transform_reduce}} &
  applies an invocable, then reduces out of order\\
\href{https://en.cppreference.com/w/cpp/algorithm/transform_exclusive_scan}{\lstinline{transform_exclusive_scan}} &
  applies an invocable, then calculates exclusive scan\\
\href{https://en.cppreference.com/w/cpp/algorithm/transform_inclusive_scan}{\lstinline{transform_inclusive_scan}} &
  applies an invocable, then calculates inclusive scan\\
\end{tabularx}
\end{frame}

\begin{frame}{Heap Algorithms}

Built-in heapsort.\\
(You must not implement the Heap in your lab using these functions, but you may use them for testing.)

\footnotesize\renewcommand{\ULdepth}{3pt}\renewcommand{\arraystretch}{1.3}
\begin{tabularx}{\textwidth}{r X}
\href{https://en.cppreference.com/w/cpp/algorithm/is_heap}{\lstinline{is_heap}} &
  checks if the given range is a max heap\\
\href{https://en.cppreference.com/w/cpp/algorithm/is_heap_until}{\lstinline{is_heap_until}} &
  finds the largest subrange that is a max heap\\
\href{https://en.cppreference.com/w/cpp/algorithm/make_heap}{\lstinline{make_heap}} &
  creates a max heap out of a range of elements\\
\href{https://en.cppreference.com/w/cpp/algorithm/push_heap}{\lstinline{push_heap}} &
  adds an element to a max heap\\
\href{https://en.cppreference.com/w/cpp/algorithm/pop_heap}{\lstinline{pop_heap}} &
  removes the largest element from a max heap\\
\href{https://en.cppreference.com/w/cpp/algorithm/sort_heap}{\lstinline{sort_heap}} &
  turns a max heap into a range of elements sorted in ascending order\\
\end{tabularx}
\end{frame}

\section{Function Objects}
\begin{frame}{Function Objects}
\begin{itemize}
\item
  Some algorithms like sort, merge, and accumulate can take a function object as an argument.
\item
  A function object is an instance of a class that overloads \lstinline{operator()}.
\item
  It is also possible to use user-written functions in place of predefined function objects.
\item
  In C++, a slang term for \textit{Function Objects} is \emph{functors}, but that term has alternate meaning.
\end{itemize}
\end{frame}

\begin{frame}[fragile]{Function Objects Example}
\begin{lstlisting}[basicstyle=\small\ttfamily]
#include <list>
#include <functional>
int main() {
  std::list<int> lst { 6, 4, 9, 1, 7 };
  lst.sort(std::greater<int>{}); // uses function object
  // for sorting in reverse order lst = { 9, 7, 6, 4, 1 }
}
\end{lstlisting}
\end{frame}

\begin{frame}[fragile]{Function Objects}
\small
The accumulate algorithm accumulates values over a range (e.g., the sum).

\begin{lstlisting}[basicstyle=\small\ttfamily]
#include <list>
#include <functional>
#include <numeric>
int main() {
  std::list<int> lst { 6, 4, 9, 1, 7 };
  int sum, fac;
  sum = std::accumulate(lst.begin(), lst.end(), 0,
    std::plus<int>{});
  sum = std::accumulate(lst.begin(), lst.end(), 0);
    // plus by default

  fac = std::accumulate(lst.begin(), lst.end(), 1,
    std::times<int>{});
}
\end{lstlisting}
\end{frame}

\begin{frame}[fragile]{User Defined Function Objects}

\begin{lstlisting}[basicstyle=\footnotesize\ttfamily]
template <typename Type>
struct SquaredSum  // user-defined function object
{
    Type operator()(Type n1, Type n2) { return n1 + n2 * n2; }
};
\end{lstlisting}
\begin{lstlisting}[basicstyle=\footnotesize\ttfamily]
#include <complex>
#include <numeric>
#include <vector>
\end{lstlisting}
\begin{lstlisting}[basicstyle=\footnotesize\ttfamily]
std::vector<std::complex<int>> vc {{2, 3}, {1, 5}, {-2, 4}};

// computes the sum of squares of a vector of complex numbers
std::complex sum_vc = std::accumulate(vc.begin(), vc.end(),
          std::complex(0, 0), SquaredSum<std::complex<int>>{});
\end{lstlisting}
\end{frame}


\begin{frame}[fragile]{User Defined Function Objects}
Unlike functions, function objects permit custom parameterization and statefullness.

\begin{lstlisting}[basicstyle=\footnotesize\ttfamily]
std::vector<double> vals1 {2, 3, 1, 5, -2, 4};
std::vector<double> vals2 {9, 8, -6, -10, -1, 8, 6};

// Square each value in both vectors and calculate the total.
CalculateAverageOfPowers squared{2}; // Function Object
squared = std::for_each(vals1.begin(), vals1.end(), squared);
squared = std::for_each(vals2.begin(), vals2.end(), squared);

std::cout << ' ' << squared.getAverage() << '\n';
\end{lstlisting}
\end{frame}

\begin{frame}[fragile]{User Defined Function Objects}
\begin{lstlisting}[basicstyle=\footnotesize\ttfamily]
class CalculateAverageOfPowers
{
public:
  CalculateAverageOfPowers(double power) 
    : total{0}, count{0}, exponent{power} {}

  void operator() (double x) {
    total += pow(x, exponent);
    ++count;
  }

  auto getAverage() const { return total / count; }

private:
  double total;
  int count;
  double exponent;
};
\end{lstlisting}
\end{frame}

\section*{Lab 17}
\begin{frame}{Lab 17: The STL and Sorting}
\centering\Large
  Let's take a look at Lab 17.
\end{frame}

\end{document}
